\section{Espacios de Hilbert y espacios de funciones}
\label{Spaces}

En esta sección se fija el marco funcional en el que se formularán los problemas inversos asociados a la tomografía computarizada. En particular, se introducen los espacios de Hilbert, los espacios $L^{p}$ y los espacios de Sobolev que se utilizarán en los capítulos posteriores. Las referencias generales son \cite{rynne_youngson2000,kirsch2011,natterer2001,kuchment2014}.


\subsection{Espacios de Hilbert}

A lo largo de este trabajo se consideran espacios vectoriales sobre $\mathbb{R}$, a menos que se indique lo contrario.

\begin{definition}
Sea $H$ un espacio vectorial real. Un \emph{producto interno} en $H$ es una aplicación
\[
\langle \cdot,\cdot \rangle : H \times H \to \mathbb{R}
\]
tal que, para todo $x,y,z \in H$ y todo $\alpha \in \mathbb{R}$, se verifican:
\begin{enumerate}
  \item $\langle \alpha x + y, z \rangle = \alpha \langle x,z \rangle + \langle y,z \rangle$;
  \item $\langle x,y \rangle = \langle y,x \rangle$ ;
  \item $\langle x,x \rangle \ge 0$;
  \item $\langle x,x \rangle = 0$ implica $x=0$ .
\end{enumerate}
En este caso se dice que $(H,\langle \cdot,\cdot \rangle)$ es un \emph{espacio con producto interno}.
\end{definition}


\begin{proposition}

    Sea $H$ un espacio con producto interno, y sean $x, y \in X$, entonces
    \begin{itemize}
        \item $|\langle x,y \rangle|^2 \leq \langle x, y \rangle \langle y, y \rangle $
        \item La función $|| \cdot || : X \to \mathbb{R}$ definida por $|| x || = \langle x, x \rangle ^{1/2}$ es una norma en $X$.
    \end{itemize}
\end{proposition}

\begin{proof}
    Pendiente
\end{proof}

\begin{definition}
Un espacio con producto interno $(H,\langle \cdot,\cdot \rangle)$ se denomina \emph{espacio de Hilbert} si es completo con respecto a la norma inducida $\|\cdot\|$, es decir, si toda sucesión de Cauchy en $(H,\|\cdot\|)$ converge en $H$.
\end{definition}

\begin{example}
Entre los ejemplos básicos de espacios de Hilbert se encuentran:
\begin{enumerate}
  \item El espacio euclídeo $\mathbb{R}^{n}$ con el producto interno
  \[
  \langle x,y \rangle = \sum_{j=1}^{n} x_{j} y_{j}.
  \]
  \item El espacio
  \[
  \ell^{2} = \Bigl\{ (x_{j})_{j \in \mathbb{N}} : \sum_{j=1}^{\infty} |x_{j}|^{2} < \infty \Bigr\},
  \]
  con producto interno
  \[
  \langle x,y \rangle = \sum_{j=1}^{\infty} x_{j} y_{j}.
  \]
  \item El espacio $L^{2}(\Omega)$, donde $\Omega \subset \mathbb{R}^{n}$ es un conjunto medible, con
  \[
  \langle f,g \rangle = \int_{\Omega} f(x)\,g(x)\,dx.
  \]
\end{enumerate}
\end{example}

En el contexto de la tomografía computarizada, los espacios de Hilbert de funciones, en particular los espacios $L^{2}(\Omega)$ y ciertos espacios de Sobolev, proporcionan un marco natural para la formulación de los problemas directo e inverso asociados a la transformada de Radon \cite{natterer2001,kuchment2014}.


\subsection{Espacios \texorpdfstring{$L^{p}$}{Lp}}

Sea $(X,\mathcal{A},\mu)$ un espacio de medida. Para $1 \le p < \infty$, se define
\[
L^{p}(X) = \left\{ f: f \text{ es medible y }  \int_{X} |f(x)|^{p}\, d\mu(x) < \infty \right\},
\]
La aplicación
\[
\|f\|_{L^{p}} := \left( \int_{X} |f(x)|^{p}\, d\mu(x) \right)^{1/p}
\]
define una norma en $L^{p}(X)$, y con esta norma $L^{p}(X)$ es un espacio de Banach \cite{rynne_youngson2000}.

Para $p = \infty$, primero definimos el supremo esencial.

\begin{definition}[Supremo esencial]
Sea $f : X \to \mathbb{R}$ una función medible. Se define el \emph{supremo esencial}
de $|f|$ como
\[
\operatorname*{ess\,sup}_{x\in X} |f(x)|
:= \inf\Bigl\{ M \ge 0 : \mu\bigl(\{x \in X : |f(x)| > M\}\bigr) = 0 \Bigr\}.
\]
\end{definition}

\begin{definition}[Espacio $L^{\infty}$]
Se define
\[
L^{\infty}(X) 
= \Bigl\{ f \text{ medible} : \operatorname*{ess\,sup}_{x\in X} |f(x)| < \infty \Bigr\},
\]
identificando funciones que coinciden casi en todas partes. La aplicación
\[
\|f\|_{L^{\infty}} := \operatorname*{ess\,sup}_{x\in X} |f(x)|
\]
define una norma en $L^{\infty}(X)$, con la cual $L^{\infty}(X)$ es un espacio de Banach.
\end{definition}

El caso $p=2$ será especialmente relevante en lo que sigue. En este caso, $L^{2}(X)$ es un espacio de Hilbert con el producto interno
\[
\langle f,g \rangle_{L^{2}} = \int_{X} f(x)\,g(x)\, d\mu(x).
\]
En particular, espacios del tipo $L^{2}(\mathbb{R}^{2})$ y $L^{2}(\Omega)$, con $\Omega \subset \mathbb{R}^{2}$, se utilizarán para modelar las imágenes y los datos de proyección en los problemas tomográficos \cite{natterer2001,kuchment2014}.


\subsection{Espacios de Sobolev}

Los espacios de Sobolev permiten cuantificar la regularidad de funciones en términos de derivadas (en sentido débil) controladas en norma $L^{2}$. Esta estructura es adecuada para formular hipótesis de suavidad sobre la solución de un problema inverso y para introducir términos de regularización.

Sea $\Omega \subset \mathbb{R}^{n}$ un abierto. Se denota por $C_{c}^{\infty}(\Omega)$ el conjunto de funciones suaves con soporte compacto contenido en $\Omega$.

\begin{definition}[Espacio $L^{1}_{\mathrm{loc}}(\Omega)$]
Se define el espacio de funciones localmente integrables como
\[
L^{1}_{\mathrm{loc}}(\Omega)
= \left\{ f : \Omega \to \mathbb{R} \text{ medible} \;\middle|\;
f \in L^{1}(K) \text{ para todo compacto } K \subset \Omega \right\}.
\]
\end{definition}


\begin{definition}[Derivada débil de primer orden]
Sea $f \in L^{1}_{\mathrm{loc}}(\Omega)$ y $1 \le j \le n$. Se dice que una función $g \in L^{1}_{\mathrm{loc}}(\Omega)$ es la \emph{derivada débil} de $f$ respecto a $x_{j}$ si
\[
\int_{\Omega} f(x)\, \frac{\partial \varphi}{\partial x_{j}}(x)\, dx
= - \int_{\Omega} g(x)\, \varphi(x)\, dx
\quad \text{para toda } \varphi \in C_{c}^{\infty}(\Omega).
\]
En este caso se escribe $g = \partial_{x_{j}} f$.
\end{definition}

\begin{definition}[Espacio de Sobolev $H^{1}(\Omega)$]
Se define
\[
H^{1}(\Omega) = \left\{ f \in L^{2}(\Omega) : \partial_{x_{j}} f \in L^{2}(\Omega) \text{ para todo } j=1,\dots,n \right\},
\]
donde las derivadas se entienden en el sentido débil. La norma en $H^{1}(\Omega)$ viene dada por
\[
\|f\|_{H^{1}(\Omega)}^{2} 
= \int_{\Omega} |f(x)|^{2}\,dx 
+ \sum_{j=1}^{n} \int_{\Omega} \left|\partial_{x_{j}} f(x)\right|^{2}\,dx.
\]
\end{definition}

\begin{proposition}
El espacio $H^{1}(\Omega)$, provisto del producto interno
\[
\langle f,g \rangle_{H^{1}(\Omega)}
= \int_{\Omega} f(x)\,g(x)\,dx
+ \sum_{j=1}^{n} \int_{\Omega} \partial_{x_{j}} f(x)\,\partial_{x_{j}} g(x)\,dx,
\]
es un espacio de Hilbert.
\end{proposition}

\begin{proof}
Pendiente
\end{proof}

Para órdenes superiores resulta conveniente utilizar la notación de multiíndices. Si $\alpha = (\alpha_{1},\dots,\alpha_{n}) \in \mathbb{N}_{0}^{n}$ y $|\alpha| = \alpha_{1} + \cdots + \alpha_{n}$, se escribe
\[
D^{\alpha} f = \partial_{x_{1}}^{\alpha_{1}} \cdots \partial_{x_{n}}^{\alpha_{n}} f,
\]
entendiendo las derivadas en sentido débil.

\begin{definition}[Espacios de Sobolev $H^{k}(\Omega)$]
Sea $k \in \mathbb{N}$. Se define
\[
H^{k}(\Omega) = \left\{ f \in L^{2}(\Omega) : D^{\alpha} f \in L^{2}(\Omega) \text{ para todo multiíndice } \alpha \text{ con } |\alpha|\le k \right\},
\]
con norma
\[
\|f\|_{H^{k}(\Omega)}^{2}
= \sum_{|\alpha|\le k} \int_{\Omega} |D^{\alpha} f(x)|^{2}\,dx.
\]
\end{definition}

Con esta norma, $H^{k}(\Omega)$ es también un espacio de Hilbert \cite{kirsch2011}. En aplicaciones a tomografía, estos espacios permiten formular términos de regularización que penalizan la falta de suavidad de la imagen a reconstruir, típicamente mediante funcionales de la forma
\[
\|Rf - g\|_{Y}^{2} + \alpha \|f\|_{H^{1}(\Omega)}^{2},
\]
donde $R$ denota la transformada de Radon, $g$ representa los datos medidos, $Y$ es un espacio de Hilbert adecuado para el sinograma y $\alpha>0$ es un parámetro de regularización. Estos aspectos se retomarán en capítulos posteriores.


